%# -*- coding:utf-8 -*-
% Course project slides for Remote Sensing Geoscience Analysis

\PassOptionsToPackage{quiet}{fontspec}
\PassOptionsToPackage{dvipsnames}{xcolor}
\documentclass[10pt,aspectratio=169]{beamer}
\usepackage{zju_beamer}
\graphicspath{{./figures/}}
\zjusectioncolor{blue}

% Lightweight helpers for compact academic slides.
\newcommand{\metric}[1]{\textcolor{MidnightBlue}{\textbf{#1}}}
\newcommand{\method}[1]{\textbf{#1}}
\newcommand{\imgfit}[2][0.95\linewidth]{%
  \IfFileExists{#2}{%
    \includegraphics[width=#1,height=0.65\textheight,keepaspectratio]{#2}%
  }{%
    \fbox{\parbox[c][0.24\textheight][c]{#1}{\centering\tiny
      图像文件暂未生成\\[0.4em]\texttt{\detokenize{#2}}\\[0.4em]
      请按 README 运行实验脚本后重新编译}}%
  }%
}

% 首页信息设置
\title[Sentinel-2 岩性聚类]{基于自编码器与 K-means 的\\Sentinel-2 多光谱遥感影像岩性聚类研究}
\subtitle{《遥感地学分析》课程大作业展示}
\author[张项翌杰]{张项翌杰}
\institute[ZJU]{浙江大学}
\date{2026}

\begin{document}

\begin{frame}
  \titlepage
\end{frame}

\begin{frame}
  \frametitle{汇报目录}
  \tableofcontents
\end{frame}

\section{研究背景}

\begin{frame}{研究动机}
  \begin{columns}[T,onlytextwidth]
    \begin{column}{0.52\linewidth}
      \begin{block}{问题背景}
        \begin{itemize}
          \item 干旱裸露矿区植被覆盖低，适合利用多光谱影像识别岩性与蚀变信息。
          \item 遥感地质填图常缺少可靠地面标签，监督分类难以直接展开。
          \item 无监督聚类可从光谱数据中发现自然分组，形成初步地质解译线索。
        \end{itemize}
      \end{block}
    \end{column}
    \begin{column}{0.43\linewidth}
      \begin{block}{核心问题}
        \small
        线性降维是否足够？\\[0.4em]
        自编码器是否能学习更适合聚类的低维光谱表示？\\[0.4em]
        聚类类别能否获得可解释的地质含义？
      \end{block}
    \end{column}
  \end{columns}
\end{frame}

\begin{frame}{研究区与数据}
  \begin{columns}[T,onlytextwidth]
    \begin{column}{0.33\linewidth}
      \begin{itemize}
        \item 研究区：新疆哈密土屋--延东斑岩型铜矿区。
        \item 数据源：Sentinel-2A MSI Level-2A 地表反射率。
        \item 影像日期：2024-08-13，Tile：T46TFM。
        \item AOI：94.30$^\circ$E--94.70$^\circ$E，42.02$^\circ$N--42.32$^\circ$N。
        \item 使用 B02--B12 中 10 个可见光、近红外和短波红外波段。
      \end{itemize}
    \end{column}
    \begin{column}{0.62\linewidth}
      \centering
      \imgfit[\linewidth]{../data/processed/preview_geology_B12_B8A_B02.png}
      \vspace{0.3em}

      \scriptsize Sentinel-2 B12/B8A/B02 地质假彩色合成
    \end{column}
  \end{columns}
\end{frame}

\section{数据与方法}

\begin{frame}{预处理流程}
  \begin{block}{从 SAFE 产品到聚类输入}
    \begin{enumerate}
      \item 解压 Sentinel-2 L2A SAFE 产品，提取 10 个光谱波段与 SCL 波段。
      \item 将所有波段统一重采样到 20 m 空间分辨率。
      \item 按 AOI 裁剪研究区，使用 SCL 去除云、阴影、水体和无效像元。
      \item 对有效像元进行光谱标准化，得到 10 维特征矩阵。
    \end{enumerate}
  \end{block}
  \begin{center}
    \Large
    最终数据集：\metric{2,772,025} 个有效像元 $\times$ \metric{10} 个光谱波段
  \end{center}
\end{frame}

\begin{frame}{四组对比方法}
  \small
  \begin{table}
    \centering
    \resizebox{\linewidth}{!}{%
      \begin{tabular}{llll}
        \hline
        方法 & 特征类型 & 维度 & 作用 \\
        \hline
        Raw K-means & 原始标准化光谱 & 10 & 无降维基线 \\
        PCA + K-means & 线性主成分 & 3 / 5 & 检验线性降维是否足够 \\
        Canonical AE + K-means & 浅层非线性隐特征 & 5 & 非线性强基线 \\
        Stacked AE + K-means & 深层非线性隐特征 & 3 / 4 / 5 & 主体方法，对齐 Naagar 等框架 \\
        \hline
      \end{tabular}}
  \end{table}
  \vspace{0.7em}
  \begin{block}{实验逻辑}
    \centering
    Raw 基线 $\rightarrow$ PCA 线性降维 $\rightarrow$ 浅层 AE 非线性表示 $\rightarrow$ 深层 SAE 层级表示
  \end{block}
\end{frame}

\begin{frame}{自编码器模型结构}
  \begin{columns}[T,onlytextwidth]
    \begin{column}{0.48\linewidth}
      \begin{block}{Canonical AE}
        \centering
        \Large
        $10 \rightarrow 16 \rightarrow 5 \rightarrow 16 \rightarrow 10$
      \end{block}
      \small
      浅层非线性变换，用于检验是否只需较简单的 encoder--decoder 即可捕获主要聚类判别结构。
    \end{column}
    \begin{column}{0.48\linewidth}
      \begin{block}{Stacked AE}
        \centering
        \Large
        $10 \rightarrow 32 \rightarrow 16 \rightarrow z \rightarrow 16 \rightarrow 32 \rightarrow 10$
      \end{block}
      \small
      逐层压缩光谱信息，实验 $z=3,4,5$，用于分析 bottleneck 维度与网络深度的影响。
    \end{column}
  \end{columns}
  \vspace{1em}
  \begin{equation*}
    \mathcal{L}_{\mathrm{MSE}}=\frac{1}{N}\sum_{i=1}^{N}\|x_i-\hat{x}_i\|^2
  \end{equation*}
\end{frame}

\begin{frame}{评价指标}
  \begin{columns}[T,onlytextwidth]
    \begin{column}{0.32\linewidth}
      \begin{block}{Silhouette $\uparrow$}
        \small
        衡量簇内紧凑性与簇间分离性，越接近 1 越好。
      \end{block}
    \end{column}
    \begin{column}{0.32\linewidth}
      \begin{block}{DBI $\downarrow$}
        \small
        衡量簇之间的平均相似度，越低表示簇间分离更清晰。
      \end{block}
    \end{column}
    \begin{column}{0.32\linewidth}
      \begin{block}{CHI $\uparrow$}
        \small
        衡量簇间离散度与簇内离散度之比，越高表示聚类结构更优。
      \end{block}
    \end{column}
  \end{columns}
  \vspace{1em}
  \begin{center}
    全量样本规模较大，三项指标在固定随机种子下抽样 \metric{100,000} 个有效像元计算。
  \end{center}
\end{frame}

\section{实验结果}

\begin{frame}{Baseline 与 PCA：线性降维提升有限}
  \begin{columns}[T,onlytextwidth]
    \begin{column}{0.33\linewidth}
      \begin{itemize}
        \item Raw K-means 通过 Elbow 方法自动选择 \metric{$k=5$}。
        \item Baseline：Silhouette=0.3709，DBI=0.8039，CHI=123,898。
        \item PCA 前 3 个主成分累计解释 \metric{99.53\%} 方差。
        \item 但 PCA $z=3$ 的 CHI 仅提升到 127,301，增益约 \metric{+2.7\%}。
      \end{itemize}
    \end{column}
    \begin{column}{0.62\linewidth}
      \centering
      \imgfit[\linewidth]{../results/cluster_pca/pca_explained_variance.png}
      \vspace{0.3em}

      \scriptsize PCA 方差解释分析
    \end{column}
  \end{columns}
\end{frame}

\begin{frame}{自编码器训练：SAE 重建更精确}
  \begin{columns}[T,onlytextwidth]
    \begin{column}{0.36\linewidth}
      \small
      \begin{table}
        \centering
        \begin{tabular}{lcc}
          \hline
          方法 & 参数量 & 最终 MSE \\
          \hline
          Canonical AE $z=5$ & 1,067 & 0.00466 \\
          SAE $z=5$ & 1,699 & \metric{0.00285} \\
          \hline
        \end{tabular}
      \end{table}
      \vspace{0.4em}
      \begin{itemize}
        \item SAE 在早期训练阶段即取得更低重建损失。
        \item 更深层级结构保留了更精细的波段间关系。
        \item 但更低 MSE 不必然意味着所有聚类指标都全面领先。
      \end{itemize}
    \end{column}
    \begin{column}{0.59\linewidth}
      \centering
      \imgfit[\linewidth]{../results/cluster_sae/sae_training_loss_z5.png}
      \vspace{0.3em}

      \scriptsize SAE $z=5$ 训练损失曲线
    \end{column}
  \end{columns}
\end{frame}

\begin{frame}{综合聚类指标对比}
  \scriptsize
  \begin{table}
    \centering
    \resizebox{\linewidth}{!}{%
      \begin{tabular}{lcccccc}
        \hline
        Method & Dim. & $k$ & Silhouette $\uparrow$ & DBI $\downarrow$ & CHI $\uparrow$ & Recon. MSE \\
        \hline
        Raw K-means & 10 & 5 & 0.3709 & 0.8039 & 123,898 & -- \\
        PCA + K-means & 3 & 5 & 0.3796 & 0.7885 & 127,301 & -- \\
        PCA + K-means & 5 & 6 & 0.3322 & 0.8864 & 112,440 & -- \\
        Canonical AE + K-means & 5 & 6 & \metric{0.4166} & 0.7338 & \metric{147,239} & 0.00466 \\
        SAE + K-means & 5 & 6 & 0.4165 & \metric{0.6981} & 144,423 & \metric{0.00285} \\
        \hline
      \end{tabular}}
  \end{table}
  \vspace{0.5em}
  \begin{itemize}
    \item 非线性自编码器特征整体优于 Raw 和 PCA 基线。
    \item Canonical AE 的 CHI 略高，说明浅层非线性已能捕获主要判别结构。
    \item SAE 的 DBI 与重建误差更优，因此作为最终地质解释主结果。
  \end{itemize}
\end{frame}

\begin{frame}{四种方法空间结果对比}
  \begin{columns}[T,onlytextwidth]
    \begin{column}{0.78\linewidth}
      \centering
      \includegraphics[width=\linewidth,height=0.80\textheight,keepaspectratio]{../results/final_figures/method_comparison_2x2.png}
    \end{column}
    \begin{column}{0.18\linewidth}
      \small
      \begin{itemize}
        \item Raw 结果存在较明显空间破碎。
        \item PCA 对整体空间格局改善有限。
        \item AE 与 SAE 产生更紧凑的低维光谱表示。
        \item SAE 在空间连续性与细粒度类别保留之间较均衡。
      \end{itemize}
    \end{column}
  \end{columns}
\end{frame}

\begin{frame}{SAE 消融：Bottleneck 维度影响聚类稳定性}
  \begin{columns}[T,onlytextwidth]
    \begin{column}{0.32\linewidth}
      \centering
      \imgfit[0.98\linewidth]{../results/cluster_sae/sae_kmeans_k5_z3.png}
      \scriptsize $z=3,k=5$：最大类 48.9\%，类别塌缩
    \end{column}
    \begin{column}{0.32\linewidth}
      \centering
      \imgfit[0.98\linewidth]{../results/cluster_sae/sae_kmeans_k6_z4.png}
      \scriptsize $z=4,k=6$：类别结构恢复
    \end{column}
    \begin{column}{0.32\linewidth}
      \centering
      \imgfit[0.98\linewidth]{../results/cluster_sae/sae_kmeans_k6_z5.png}
      \scriptsize \metric{$z=5,k=6$}：最终解释配置
    \end{column}
  \end{columns}
  \vspace{0.5em}
  \begin{center}
    过窄 bottleneck 会压缩掉关键光谱差异，$z=5$ 在压缩与可分性之间更稳健。
  \end{center}
\end{frame}

\section{地质解释}

\begin{frame}{聚类类别的光谱响应}
  \begin{columns}[T,onlytextwidth]
    \begin{column}{0.70\linewidth}
      \centering
      \imgfit[\linewidth]{../results/cluster_spectra/cluster_spectral_profiles.png}
    \end{column}
    \begin{column}{0.26\linewidth}
      \small
      \begin{block}{Cluster 1}
        \begin{itemize}
          \item B11/B12 = \metric{0.992}，为 6 类中最低。
          \item B11 $\rightarrow$ B12 出现明显反射率下降。
          \item B12 约 2190 nm，对 Al--OH、Mg--OH 蚀变矿物吸收敏感。
        \end{itemize}
      \end{block}
      \small
      因此 Cluster 1 可解释为高置信度热液蚀变候选类。
    \end{column}
  \end{columns}
\end{frame}

\begin{frame}{最终遥感地质解释图}
  \begin{columns}[T,onlytextwidth]
    \begin{column}{0.78\linewidth}
      \centering
      \includegraphics[width=\linewidth,height=0.72\textheight,keepaspectratio]{../results/final_figures/python_final_interpretation_map.png}
    \end{column}
    \begin{column}{0.18\linewidth}
      \small
      \begin{itemize}
        \item 底图：B12/B8A/B02 地质假彩色影像。
        \item 叠加层：SAE $z=5,k=6$ 聚类结果。
        \item Cluster 1 以轮廓强调，作为蚀变候选区域。
      \end{itemize}
      \vspace{0.2em}
      \begin{block}{注意}
        \tiny
        该图是基于无监督聚类的光谱推断，未经过野外验证，不应视为正式地质图。
      \end{block}
    \end{column}
  \end{columns}
\end{frame}

\section{结论}

\begin{frame}{结论与展望}
  \small
  \begin{block}{主要结论}
    \begin{enumerate}
      \item 构建了土屋--延东地区 Sentinel-2 10 波段预处理与无监督岩性聚类流程。
      \item 非线性 AE 表示整体优于 Raw 与 PCA；SAE 在 DBI 和 MSE 上更优，适合作为最终解释主结果。
      \item Cluster 1 呈现明确 SWIR 吸收特征，可作为热液蚀变候选类。
    \end{enumerate}
  \end{block}
  \begin{block}{局限与后续工作}
    \scriptsize
    局限：缺少地面真值验证；全连接 AE 未利用空间邻域；仅使用单时相影像；K-means 对非球形光谱簇适应性有限。后续可引入野外样点、卷积自编码器、多时相数据和更灵活的聚类模型。
  \end{block}
\end{frame}

\section{参考文献}
\begin{frame}[allowframebreaks]
  \frametitle{参考文献}
  \reference
\end{frame}

\end{document}
